% Part: normal-modal-logic
% Chapter: frame-correspondence
% Section: frames

\documentclass[../../../include/open-logic-section]{subfiles}

\begin{document}

\olfileid{nml}{frd}{fra}

\olsection{框架}

\begin{defn}
  \emph{框架}是一个有序对 $\mModel{F} = \tuple{W,R}$，其中 $W$ 是一个非空世界集合，而 $R$ 是 $W$ 上的一个二元关系。模型 $\mModel{M}$ \emph{基于}框架 $\mModel{F} = \tuple{W,R}$，当且仅当存在某个赋值 $V$ 使得 $\mModel{M} = \tuple{W, R, V}$。
\end{defn}

\begin{defn}
  若 $\mModel{F}$ 是框架，如果每一个基于 $\mModel{F}$ 的模型 $\mModel{M}$ 都满足 $\mSat{M}{!A}$，则称 $!A$ \emph{在 $\mModel{F}$ 中有效}，记作 $\mModel{F} \Entails !A$。

  若 $\mClass{F}$ 是一个框架类，当且仅当对每个框架 $\mModel{F} \in \mClass{F}$ 都有 $\mModel{F} \Entails !A$ 时，我们称 $!A$ \emph{在 $\mClass{F}$ 中有效}，记作 $\mClass{F} \Entails !A$。
\end{defn}

框架之所以有趣，是因为模式与可达关系~$R$ 的性质之间的对应关系是在框架层面而非\emph{模型层面}成立的。例如，尽管公理 \Ax{T} 在所有自反模型中为真，但并非 \Ax{T} 在其中为真的每个模型都是自反的。然而，\emph{确实}不仅 \Ax{T} 在所有自反\emph{框架}上都是\emph{有效的}，而且 \Ax{T} 在其上有效的每个框架也都是自反的。

\begin{rem}
框架类中的有效性是模型类中有效性概念的一个特例：$\mClass{F} \Entails !A$ 当且仅当 $\mClass{C} \Entails !A$，其中 $\mClass{C}$ 是由所有基于 $\mClass{F}$ 中某个框架的模型构成的类。

显然，如果一个公式或模式是有效的，即相对于\emph{所有}模型类有效，那么它相对于任何框架类 $\mClass{F}$ 也是有效的。
\end{rem}

\end{document}